\documentclass[11pt,a4paper]{article}
\usepackage[utf8]{inputenc}
\usepackage{geometry}
\usepackage{booktabs}
\usepackage{enumitem}
\usepackage{titlesec}
\usepackage{xcolor}
\usepackage{hyperref}
\usepackage{graphicx}
\usepackage{tabularx}
\usepackage{listings}
\usepackage{tcolorbox}

\tcbuselibrary{skins}
\geometry{margin=1in}

\titleformat{\section}{\large\bfseries\color{blue!70!black}}{}{0em}{}[\titlerule]
\titleformat{\subsection}{\normalsize\bfseries\color{blue!50!black}}{}{0em}{}
\newcommand{\fieldname}[1]{\textcolor{red!70!black}{\texttt{#1}}}
\newcommand{\notesquare}{\textcolor{gray!70!black}{\rule{0.45em}{0.45em}}}
\newcommand{\jsonkey}[1]{\textcolor{blue!70!black}{\texttt{#1}}}
\newcommand{\jsonstring}[1]{\textcolor{green!45!black}{\texttt{#1}}}
\newcommand{\jsonnumber}[1]{\textcolor{orange!85!black}{\texttt{#1}}}
\newcommand{\jsoncolon}{\textcolor{red!70!black}{\texttt{:}}}

\lstset{
  backgroundcolor=\color{gray!10},
  frame=none,
  xleftmargin=1em,
  xrightmargin=1em,
  aboveskip=0.2em,
  belowskip=0em,
  columns=fullflexible,
  keepspaces=true,
  showstringspaces=false,
  breaklines=true,
  breakatwhitespace=true
}

\lstdefinestyle{pageindexjson}{
  numbers=none,
  linewidth=\textwidth,
  xleftmargin=-\leftmargini,
  xrightmargin=0em,
  basicstyle=\ttfamily\footnotesize\color{gray!80!black},
  escapeinside={(*@}{@*)}
}

\lstdefinestyle{pageTextBlock}{
  numbers=none,
  breaklines=true,
  linewidth=\textwidth,
  xleftmargin=-\leftmargini,
  xrightmargin=0em,
  basicstyle=\ttfamily\footnotesize\color{gray!80!black}
}

\begin{document}

\section{PageIndex Implementation Flow}


\begin{lstlisting}[
    numbers=none,
    breaklines=true,
    linewidth=\textwidth,
    xleftmargin=-\leftmargini,
    xrightmargin=0em,
    basicstyle=\ttfamily\footnotesize\color{gray!70!black},
    mathescape=true,
    literate={->}{$\rightarrow$}{2}
]
document -> page text -> ToC detection -> tree construction -> index verification -> retrieval-ready document record
\end{lstlisting}

\begin{enumerate}
    \item \textbf{Parse pages}: Extract plain text from the pdf.
    \item \textbf{Detect existing ToC}: Decide whether the document already contain a ToC.
    \item \textbf{Choose tree-building path}: 
    \begin{itemize}
        \item \textbf{ToC w/ page numbers}: Converts to \textbf{JSON} with page numbers. 
        \item \textbf{ToC w/o page numbers}: LLM identifies where each section starts.
        \item \textbf{If no ToC exists}: LLM infers the hierarchy directly from page-marked text.
    \end{itemize}   
    \item \textbf{Group pages}: Long documents are split into token-safe groups before prompting. This keeps each LLM call within context limits and avoids forcing the model to read the whole document at once.
    \item \textbf{Verify and repair indexes}: Generated section boundaries are checked against the document text. If a section appears to start on the wrong page, the system searches nearby pages and asks the LLM to repair the boundary.
    \item \textbf{Convert section list to tree}: The cleaned section list is converted into a recursive hierarchy. Each node receives a \fieldname{start\_index}, an \fieldname{end\_index}, and child nodes when the section contains smaller sections.
    \item \textbf{Split large nodes}: Sections that are too large are processed again at a finer level, creating deeper subtrees only where the document needs more detail.
    \item \textbf{Add retrieval fields}: The final tree can include \fieldname{node\_id}, raw \fieldname{text}, node \fieldname{summary}, and a document-level description depending on the indexing settings.
    \item \textbf{Store retrieval-ready data}: The system stores the document record, tree, page count, extracted page text, and optional descriptions so later retrieval can return either the compact tree or the raw evidence.
    \item \textbf{Support query-time retrieval}: At query time, the system can expose document metadata, a tree without raw text, and selected raw page evidence for answer generation.
\end{enumerate}

\noindent\textbf{LLM usage note:} These steps can use the same model, but they are separate LLM calls. The pipeline passes explicit data between steps, such as extracted page text, ToC JSON, selected \fieldname{doc\_id}, selected \fieldname{node\_id}, and retrieved evidence. It does not rely on one shared hidden chat history.

\newpage
\section{Indexing Step 1: Create PageIndex Tree for Each Document Using LLM }
\begin{itemize}
    \item \textbf{Input}: Document
    \item \textbf{Output}: \textbf{JSON-based hierarchical structure} that represents \textbf{ToC}
    \item \textbf{LLM role}: The LLM can be used to identify section boundaries, generate titles, write summaries, and decide how detailed the tree should be.
\end{itemize}

\noindent\textbf{Implementation detail:}
\begin{enumerate}
    \item \textbf{Extract page text}: Extract ordered page-level text from document.
    \begin{itemize}
        \item The actual extraction output is a list of \texttt{(page\_text, token\_count)} pairs.
        \item The physical page index is implied by the list position: the first pair is page 1, the second pair is page 2, etc.

    \begin{lstlisting}[style=pageTextBlock]
[
  ("Annual Report 2023 ...", 742),
  ("Financial Stability ...", 691)
]
\end{lstlisting}
        \item Before the LLM receives the page text, the ordered pages are wrapped with physical page tags:
  
\begin{lstlisting}[style=pageTextBlock]
<physical_index_1>
Annual Report 2023 ...
<physical_index_1>

<physical_index_2>
Financial Stability ...
<physical_index_2>
\end{lstlisting}

\end{itemize} 

    \item \textbf{Detect existing ToC using LLM}: 
    \begin{itemize}
        \item Using LLM to label whether each page is a ToC. Since existing structure usually reflects hierarchy structure.
\begin{lstlisting}[style=pageTextBlock]
toc_page_list = [1, 2]
\end{lstlisting}
        \item Page-check behavior:
        
        \begin{tabularx}{\linewidth}{@{}p{0.30\linewidth}X@{}}
        \toprule
        \textbf{Situation} & \textbf{Pages checked} \\
        \midrule
        \textbf{No ToC found} & Up to \textbf{\textasciitilde20 pages} (PDF pages 1--20), then stop \\
        \textbf{ToC found early} & Only until the ToC block ends (could be 3 pages, could be 8) \\
        \textbf{ToC spans 20 pages} & \textbf{More than 20} -- keeps going while still inside the ToC block \\
        \bottomrule
        \end{tabularx}
        
\end{itemize} 

    \item \textbf{Detect existing ToC using LLM}: 
        \item Concatenate those pages' text and feed to LLM to detect whether page numbers are given for sections in the ToC.
        \begin{lstlisting}[style=pageTextBlock]
Example input to that one call:

Contents
1. Introduction .............. 3
2. Monetary Policy ........... 9
3. Financial Stability ...... 21

-> page_index_given_in_toc: "yes"

vs.

Contents
1. Introduction
2. Monetary Policy
3. Financial Stability

-> page_index_given_in_toc: "no"
\end{lstlisting}
    
        
    \item \textbf{Build the first tree}: If a ToC exists, the system normalizes it into JSON and aligns its section pages to physical pages. If no ToC exists, the LLM reads page-marked text and proposes section titles, hierarchy, and start pages.
    \item \textbf{Handle long documents}: The document can be divided into token-safe groups. The LLM reads one group at a time and returns partial structure, then the system merges the partial structure into one document tree.
    \item \textbf{Verify page boundaries}: Section titles and start pages are checked against the extracted text. When a boundary is wrong, the system searches nearby pages and repairs the index.
    \item \textbf{Create recursive nodes}: The final structure is converted into nested nodes. Parent nodes cover broad sections; child nodes cover more specific subsections.
    \item \textbf{Control detail level}: The tree becomes deeper when a section is too long or too dense. Short sections can remain as leaf nodes to avoid unnecessary complexity.
    \item \textbf{Attach retrieval fields}: Each node receives a stable \fieldname{node\_id}, page range, optional summary, and optional raw text pointer.
\end{enumerate}

\begin{itemize}
    \item \textbf{Nodes} represent a section (e.g. chapter, paragraph, page)
\begin{lstlisting}[
    language=Java,
    numbers=none,
    breaklines=false,
    linewidth=\textwidth,
    xleftmargin=-\leftmargini,
    xrightmargin=0em,
    morekeywords={Node,string,object},
    keywordstyle=\color{gray!70!black},
    morekeywords={[2]node_id,title,description,metadata,sub_nodes},
    keywordstyle={[2]\color{blue}},
    commentstyle=\color{gray!70!black}\itshape,
    basicstyle=\ttfamily\normalsize\color{gray!70!black},
    literate={:}{{\textcolor{red}{:}}}{1}
]
Node {
  node_id: string,        // Unique node identifier
  title: string,           // Human-readable label or title
  description: string,    // Optional detailed explanation of the node
  metadata: object,       // Arbitrary key-value pairs for context or attributes
  sub_nodes: [Node]       // Array of child nodes (recursive structure)
}
\end{lstlisting} 
    \begin{itemize}[label=\notesquare]
        \item \fieldname{node\_id}: Serves as a reference key to locate the corresponding raw data.
        \item \fieldname{sub\_nodes}: Field allows recursive nesting, forming a full ToC tree.
        \item \fieldname{metadata}: Field can store semantic information (e.g. document type, author, timestamp, or relevance scores.)
    \end{itemize}
    
    \item \textbf{\large ToC Example:}
    \begin{lstlisting}[style=pageindexjson,basicstyle=\ttfamily\normalsize\color{gray!80!black}]
...
{
  (*@\jsonkey{"node\_id"}@*)(*@\jsoncolon@*) (*@\jsonstring{"0006"}@*),
  (*@\jsonkey{"title"}@*)(*@\jsoncolon@*) (*@\jsonstring{"Financial Stability"}@*),
  (*@\jsonkey{"start\_index"}@*)(*@\jsoncolon@*) (*@\jsonnumber{21}@*),
  (*@\jsonkey{"end\_index"}@*)(*@\jsoncolon@*) (*@\jsonnumber{22}@*),
  (*@\jsonkey{"summary"}@*)(*@\jsoncolon@*) (*@\jsonstring{"The Federal Reserve ..."}@*),
  (*@\jsonkey{"sub\_nodes"}@*)(*@\jsoncolon@*) [
    {
      (*@\jsonkey{"node\_id"}@*)(*@\jsoncolon@*) (*@\jsonstring{"0007"}@*),
      (*@\jsonkey{"title"}@*)(*@\jsoncolon@*) (*@\jsonstring{"Monitoring Financial Vulnerabilities"}@*),
      (*@\jsonkey{"start\_index"}@*)(*@\jsoncolon@*) (*@\jsonnumber{22}@*),
      (*@\jsonkey{"end\_index"}@*)(*@\jsoncolon@*) (*@\jsonnumber{28}@*),
      (*@\jsonkey{"summary"}@*)(*@\jsoncolon@*) (*@\jsonstring{"The Federal Reserve's monitoring ..."}@*)
    },
    {
      (*@\jsonkey{"node\_id"}@*)(*@\jsoncolon@*) (*@\jsonstring{"0008"}@*),
      (*@\jsonkey{"title"}@*)(*@\jsoncolon@*) (*@\jsonstring{"Domestic and International Cooperation and Coordination"}@*),
      (*@\jsonkey{"start\_index"}@*)(*@\jsoncolon@*) (*@\jsonnumber{28}@*),
      (*@\jsonkey{"end\_index"}@*)(*@\jsoncolon@*) (*@\jsonnumber{31}@*),
      (*@\jsonkey{"summary"}@*)(*@\jsoncolon@*) (*@\jsonstring{"In 2023, the Federal Reserve collaborated ..."}@*)
    }
  ]
}
...
\end{lstlisting}
\begin{itemize}[label=\notesquare]
    \item \fieldname{start\_index / end\_index}: Start page / End page
    \item Each node in ToC links to corresponding \textbf{raw content} (e.g. text, images, tables)
    \begin{lstlisting}[
    language=Java,
    numbers=none,
    breaklines=false,
    linewidth=\textwidth,
    xleftmargin=-\leftmargini,
    xrightmargin=0em,
    morekeywords={Node,string,object},
    keywordstyle=\color{gray!70!black},
    morekeywords={[2]node_id,title,description,metadata,sub_nodes},
    keywordstyle={[2]\color{blue}},
    commentstyle=\color{gray!70!black}\itshape,
    basicstyle=\ttfamily\normalsize\color{gray!70!black},
    literate={:}{{\textcolor{red}{:}}}{1},
    mathescape=true,                    
    literate={->}{$\rightarrow$}{2}     
]
    node_id -> node_content (raw content, extracted text, images, etc.)
\end{lstlisting} 
\end{itemize}
    \begin{itemize} 
        \item This mapping enables the LLM to \textbf{select and retrieve} specific nodes as needed.
        \item \textbf{JSON-based ToC index} resides within the LLM's active reasoning context. 
    \end{itemize}
\end{itemize}


\section{Indexing Step 2: Create Documents Record Using LLM (Can be Same LLM)}

The document-level record answers: ``Which document should answer the query?'' This record should be small, stable, and easy to filter.

\begin{itemize}
    \item \textbf{Main fields}: \fieldname{doc\_id}, file name, title, source path, document type, author, date, language, and access permission.
    \item \textbf{Description}: A short document-level summary used for document search.
    \begin{itemize}
        \item For a small collection, create the PageIndex ToC tree first, then derive the document-level description from the root summary, top-level node summaries, and metadata.
        \item For a large collection, create a cheaper document-level record first using metadata, filename, abstract, title page, or a short extracted preview. Generate full PageIndex trees only for documents that are likely to be searched.
    \end{itemize}
    \item \textbf{Storage pointer}: The record should point to where the full ToC tree, raw text, page images, or original PDF are stored.
\end{itemize}

\noindent\textbf{Implementation detail:}
\begin{enumerate}
    \item \textbf{Create a stable document id}: Each document receives an id that can be used throughout the whole pipeline. Query-time retrieval returns this id before searching inside the document tree.
    \item \textbf{Collect basic metadata}: The record stores stable facts such as file name, title, source, year, type, language, and permission. These fields are useful for filtering and access control.
    \item \textbf{Generate or derive a description}: The description is a short document-level summary. It can be written by the LLM from the root and top-level tree summaries, or from a short preview when the full tree has not been created yet.
    \item \textbf{Store pointers}: The record keeps locations for the original file, tree, extracted text, page images, or other raw evidence. The record itself stays small while pointing to larger stored objects.
    \item \textbf{Separate document record from tree}: The document record is used to decide \fieldname{which document} to search. The PageIndex tree is used later to decide \fieldname{which section} inside that document to search.
    \item \textbf{Keep the record updateable}: If the document is re-indexed, the record can keep the same \fieldname{doc\_id} while updating its description, metadata, and storage pointers.
\end{enumerate}

\noindent\textbf{\large Document-Level Record Example:}
\begin{lstlisting}[style=pageindexjson,breaklines=false,basicstyle=\ttfamily\normalsize\color{gray!80!black}]
{
  (*@\jsonkey{"doc\_id"}@*)(*@\jsoncolon@*) (*@\jsonstring{"fed\_annual\_2023"}@*),
  (*@\jsonkey{"file\_name"}@*)(*@\jsoncolon@*) (*@\jsonstring{"annual-report-2023.pdf"}@*),
  (*@\jsonkey{"title"}@*)(*@\jsoncolon@*) (*@\jsonstring{"Federal Reserve Annual Report 2023"}@*),
  (*@\jsonkey{"description"}@*)(*@\jsoncolon@*) (*@\jsonstring{"Annual report for document-level search."}@*),
  (*@\jsonkey{"metadata"}@*)(*@\jsoncolon@*) {
    (*@\jsonkey{"source"}@*)(*@\jsoncolon@*) (*@\jsonstring{"Federal Reserve"}@*),
    (*@\jsonkey{"year"}@*)(*@\jsoncolon@*) (*@\jsonnumber{2023}@*),
    (*@\jsonkey{"document\_type"}@*)(*@\jsoncolon@*) (*@\jsonstring{"annual\_report"}@*),
    (*@\jsonkey{"language"}@*)(*@\jsoncolon@*) (*@\jsonstring{"en"}@*)
  },
  (*@\jsonkey{"toc\_tree\_uri"}@*)(*@\jsoncolon@*) (*@\jsonstring{"storage/trees/fed\_annual\_2023.json"}@*),
  (*@\jsonkey{"raw\_file\_uri"}@*)(*@\jsoncolon@*) (*@\jsonstring{"storage/files/annual-report-2023.pdf"}@*)
}
\end{lstlisting}

\begin{lstlisting}[
    numbers=none,
    breaklines=false,
    linewidth=\textwidth,
    xleftmargin=-\leftmargini,
    xrightmargin=0em,
    basicstyle=\ttfamily\normalsize\color{gray!70!black},
    mathescape=true,
    literate={->}{$\rightarrow$}{2}
]
doc_id -> document record -> ToC tree / node_map / raw files
\end{lstlisting}

\section{Indexing Step 3: Build Optional Document-Search Index}

This step is needed when the system has many documents. It helps choose which documents should be searched before running tree search inside them.

\begin{itemize}
    \item \textbf{SQL metadata search}: Good for exact filters such as year, source, company, document type, or permission.
    \item \textbf{Vector search}: Good for semantic matching against document descriptions, titles, abstracts, or top-level summaries.
\end{itemize}

\noindent\textbf{Implementation detail:}
\begin{enumerate}
    \item \textbf{Choose searchable fields}: The index can use title, description, metadata, abstract, root summary, and top-level section summaries. Raw full-document text is usually too large for document-level search.
    \item \textbf{Build metadata filters}: Structured fields are stored in a database so the system can apply exact filters before semantic search. Examples include year, source, language, document type, and permission.
    \item \textbf{Build semantic search}: Document descriptions and high-level summaries can be embedded into vectors. The user query is embedded at query time and matched against these document vectors.
    \item \textbf{Combine retrieval signals}: Production systems often combine metadata filters, keyword search, vector similarity, and optional LLM reranking. This reduces the number of documents passed into tree search.
    \item \textbf{Return candidate documents}: The output is a small ranked set of \fieldname{doc\_id} values. Tree search then runs only inside those selected documents.
    \item \textbf{Handle scale}: For small collections, this step can be simple or skipped. For large collections, it is important because generating and searching every full tree for every query is expensive.
    \item \textbf{Keep the index synchronized}: When documents are added, deleted, or re-indexed, the document-search index should update together with the document records.
\end{enumerate}

\section{Query Step 1: Document Search (Can be same LLM)}

It uses the document-level index and returns candidate \fieldname{doc\_id} values.

\begin{itemize}
    \item \textbf{Input}: User question, optional filters, and conversation context if needed.
    \item \textbf{Output}: A ranked list of candidate documents.
\end{itemize}

\section{Query Step 2: Tree Search (Can be same LLM)}

The LLM reads the PageIndex ToC tree for each selected document, usually without the full raw document text, and selects relevant \fieldname{node\_id} values.

\begin{itemize}
    \item \textbf{Input}: User question, selected \fieldname{doc\_id}, and that document's ToC tree.
    \item \textbf{Output}: Relevant \fieldname{node\_id} values, often with short reasons.
\end{itemize}

\section{Query Step 3: Retrieve Raw Content}

After tree search selects nodes, the system uses \fieldname{node\_map} to fetch the real evidence.

\section{Query Step 4: Answer the Question (C)}

\begin{itemize}
    \item \textbf{Input}: User question, retrieved evidence, page numbers, \fieldname{doc\_id}, and \fieldname{node\_id}.
    \item \textbf{Output}: A grounded answer with citations to pages, nodes, or documents.
\end{itemize}


\section{Different Methods}

\subsection{Document Search vs. Tree Search}

In a multi-document system, retrieval is usually split into two levels.

\begin{itemize}
    \item \textbf{Document search}: Find the relevant document records. The output is one or more \fieldname{doc\_id} values.
    \item \textbf{Tree search}: Search inside each selected document's PageIndex ToC tree. The output is one or more \fieldname{node\_id} values.
\end{itemize}

\begin{lstlisting}[
    numbers=none,
    breaklines=false,
    linewidth=\textwidth,
    xleftmargin=-\leftmargini,
    xrightmargin=0em,
    basicstyle=\ttfamily\normalsize\color{gray!70!black},
    mathescape=true,
    literate={->}{$\rightarrow$}{2}
]
query -> doc-search -> doc_id -> tree-search -> node_id -> raw content -> answer
\end{lstlisting}

\begin{itemize}
    \item \textbf{Document search tutorial} (\texttt{doc-search}): Shows how to choose documents using metadata, semantics, or LLM-generated descriptions.
    \item \textbf{Tree search tutorial} (\texttt{tree-search}): Shows how to choose relevant nodes inside one known document.
\end{itemize}

\subsection{Manual Text RAG: \texttt{demo/PageIndex\_Simple}}

This demo starts at the \textbf{tree-search} level because the document is already known.

\begin{enumerate}
    \item \textbf{Convert PDF to Tree(ToC)}: PageIndex creates node titles, summaries, page indexes, raw text, and node ids.
    \item \textbf{Tree search}: Give the LLM the question and the tree without raw text. The LLM returns relevant \fieldname{node\_id} values.
    \item \textbf{Map nodes to content}: Use \fieldname{node\_map} to locate the full content for each selected id.
    \item \textbf{Generate answer}: Feed only the selected node text to the LLM to answer.
\end{enumerate}

\noindent\textbf{Production note:} For many documents, add a document-search step before this flow.

\subsection{Vision RAG: \texttt{demo/Vision\_PageIndex}}

This demo also starts from one known document, but the final answer is generated from rendered page images.

\begin{enumerate}
    \item \textbf{Convert PDF to Tree(ToC)}: PageIndex creates node titles, summaries, page indexes, raw text, and node ids.
    \item \textbf{Render pages}: Convert PDF pages into image files.
    \item \textbf{Tree search}: Use a multimodal model as a text-only model to select nodes from the tree.
    \item \textbf{Collect page images}: Convert selected nodes into page ranges and get page images.
    \item \textbf{Generate visual answer}: Send the question and selected page images to a vision-language model.
\end{enumerate}

\noindent\textbf{Production note:} Use this when figures, tables, scanned pages, or visual layout are important.

\subsection{Chat Quickstart: \texttt{demo/Chat\_Quickstart\_PageIndex}}

This demo hides the tree-search and answer-generation details behind the PageIndex chat API.

\begin{enumerate}
    \item \textbf{Submit document}: Upload a PDF and keep the returned \fieldname{doc\_id}.
    \item \textbf{Check status}: Wait until PageIndex finishes processing.
    \item \textbf{Ask question}: Call \texttt{pi\_client.chat\_completions(...)} with the user question and \fieldname{doc\_id}.
    \item \textbf{Stream answer}: PageIndex handles retrieval and answer generation internally.
\end{enumerate}

\subsection{Agentic Retrieval: \texttt{demo/Agentic\_PageIndex}}

Agentic retrieval is not just ``LLM retrieves then LLM answers.'' Normal PageIndex can also do that. The difference is that retrieval becomes a \textbf{tool} inside a larger agent workflow.

\begin{enumerate}
    \item \textbf{Submit document}: Upload a PDF and keep \fieldname{doc\_id}.
    \item \textbf{Ask retrieval prompt}: Request raw relevant evidence in a structured JSON format.
    \item \textbf{Parse JSON}: Extract page numbers and evidence content from the model output.
    \item \textbf{Check / reuse evidence}: An agent can validate, compare, call more tools, retrieve again, or pass evidence to another step.
    \item \textbf{Generate final response if needed}: The final answer may be produced after the evidence is collected.
\end{enumerate}

\noindent\textbf{Key distinction:}
\begin{itemize}
    \item \textbf{Normal PageIndex}: Retrieval is a step inside a fixed QA pipeline.
    \item \textbf{Agentic PageIndex}: Retrieval is an iterative tool that an agent can call conditionally.
\end{itemize}

\subsection{Comparison}

\begin{tabularx}{\textwidth}{@{}p{0.22\textwidth}p{0.30\textwidth}p{0.38\textwidth}@{}}
\toprule
\textbf{Method} & \textbf{Main Role} & \textbf{Flow Summary} \\
\midrule
Document Search & Select document(s) & Query $\rightarrow$ document index $\rightarrow$ \fieldname{doc\_id} \\
Tree Search & Select section(s) & \fieldname{doc\_id} $\rightarrow$ ToC tree $\rightarrow$ \fieldname{node\_id} \\
Manual Text RAG & Custom text QA & Tree $\rightarrow$ node ids $\rightarrow$ node text $\rightarrow$ answer \\
Vision RAG & Visual document QA & Tree $\rightarrow$ node ids $\rightarrow$ page images $\rightarrow$ VLM answer \\
Chat Quickstart & Simple hosted QA & Document id $\rightarrow$ PageIndex chat API $\rightarrow$ streamed answer \\
Agentic Retrieval & Structured evidence tool & Query $\rightarrow$ retrieval prompt $\rightarrow$ JSON evidence $\rightarrow$ agent use \\
\bottomrule
\end{tabularx}

\end{document}
